В качестве альтернативного пути профессионального развития рассматривается самозанятость в виде фриланса в сфере оказания услуг по разработке программного обеспечения, созданию веб-инструментов и автоматизации бизнес-процессов. 

Особенности режима самозанятости предполагают лимит годового дохода до $2,4$ млн руб., самостоятельный поиск клиентов и отсутствие социальных гарантий, однако предоставляют гибкость рабочего графика и возможность добровольного пенсионного страхования.

\subsection{Описание клиента и продвижение услуг}

Для успешного позиционирования на рынке фриланса сформирован профиль потенциального клиента по методу «карта эмпатии» (\cref{tab:empathy_map}). Целевая аудитория — индивидуальные предприниматели и представители малого бизнеса, нуждающиеся в цифровизации и автоматизации рутинных задач.

\begin{longtblr}[
    caption = {Аватар клиента},
    label = {tab:empathy_map},
]{
    colspec = {X[1,j,m] X[1,j,m]},
    hlines, vlines,
    rowhead = 0,
}
    \SetCell[c=2]{c, bg=gray9} \textit{Аватар: Владелец малого бизнеса / Руководитель IT-проекта} \\
    \textit{Что думает и чувствует?} & \textit{Что видит (информация из внешней среды)?} \\
    Опасается срыва сроков, некачественного кода и исчезновения исполнителя. Желает оптимизировать расходы и автоматизировать рутину для роста прибыли. Успех измеряется стабильно работающим продуктом. & 
    Видит конкурентов, внедряющих IT-решения. Сталкивается с большим количеством некомпетентных фрилансеров на биржах. Изучает кейсы успешной автоматизации в своей нише. \\
    \textit{Что слышит (информация из внешней среды)?} & \textit{Что говорит и делает? Как выглядит?} \\
    Слышит рекомендации коллег по выбору надежных подрядчиков. Получает информацию о новых трендах (боты, парсинг, ИИ) из бизнес-подкастов и профильных Telegram-каналов. & 
    Запрашивает портфолио и реальные кейсы. Требует прозрачных смет и гарантий. Общается преимущественно в мессенджерах (Telegram) или через Zoom. Выглядит прагматичным и ориентированным на результат. \\
\end{longtblr}

В качестве основных каналов продвижения услуг выбраны:
\begin{itemize}
    \item специализированные фриланс-площадки (Хабр Фриланс, FL.ru);
    \item профессиональные сообщества и чаты в мессенджере Telegram;
    \item публикация экспертного контента (видеодемонстрации работы скриптов и алгоритмов) на YouTube.
\end{itemize}

Выбор данных каналов обусловлен их высокой концентрацией целевой аудитории и возможностью бесплатного роста. На начальном этапе расходы на платное продвижение не планируются, бюджет составляет $0$ руб. в год.

\subsection{Анализ рынка и формулировка предложения}

Для оценки конкурентной среды проведен анализ аналогичных предложений на рынке IT-фриланса (\cref{tab:freelance_market}).

\begin{longtblr}[
    caption = {Анализ рынка фриланс},
    label = {tab:freelance_market},
]{
    colspec = {X[1.2,l,m] X[1,l,m] X[1,l,m] X[1,l,m] X[1,l,m] X[1,l,m]},
    hlines, vlines,
    row{1} = {bg=gray9, c},
    rowhead = 1,
}
    \textit{Характеристика} & \textit{Предложение 1} & \textit{Предложение 2} & \textit{Предложение 3} & \textit{Предложение 4} & \textit{Предложение 5} \\
    \textit{Профиль} & Разработка Telegram-ботов & Создание парсеров данных & Fullstack веб-разработка & Разработка баз данных & Автоматизация процессов в ОС \\
    \textit{Знания} & API, алгоритмы & HTML, веб & Базы данных, веб & SQL, архитектура БД & ОС, скриптинг \\
    \textit{Умения} & Интеграция платежей, памяти & Сбор данных & Верстка, бэкенд & Оптимизация запросов & Написание макросов \\
    \textit{ПО / Стек} & Python, Aiogram & Python & JS, React, Node.js & C/C++, PostgreSQL & Bash, Python, Lua \\
    \textit{Иностранные языки} & Английский (A2) & Английский (A2) & Английский (B1) & Английский (A2) & Не требуется \\
    \textit{Личные особенности} & Всегда на связи & Внимательность к данным & Соблюдение сроков & Скурпулезность и качество & Скорость работы \\
    \textit{Требуемая плата, р./час} & $1000$ & $800$ & $1500$ & $1200$ & $600$ \\
\end{longtblr}

\textit{Формулировка предложения:} Оказание услуг по разработке программного обеспечения, созданию систем управления базами данных (СУБД), скриптов автоматизации и Telegram-ботов. Гарантия соблюдения сроков и чистоты кода. Стоимость услуг: от $600$ руб./час.

\subsection{Анализ разрыва и план по его устранению}

Сравнение текущих компетенций (\cref{tab:swot_analysis}) с требованиями рынка фриланса выявляет дефицит преимущественно в области мягких навыков --- навыки продаж, юридическая грамотность, самопрезентация --- и некоторых специфических веб-технологиях. Дополнительный индивидуальный план развития представлен в \cref{tab:idp_freelance}.

\begin{longtblr}[
    caption = {Дополнительный индивидуальный план развития (фриланс)},
    label = {tab:idp_freelance},
]{
    colspec = {X[0.8,l,m] X[1.5,l,m] X[1.5,l,m] X[1,l,m] X[1,l,m] X[0.8,l,m]},
    hlines, vlines,
    row{1} = {bg=gray9, c},
    rowhead = 1,
}
    \textit{Зона развития} & \textit{Конкретная цель} & \textit{Мероприятия по достижению} & \textit{Ключевой показатель} & \textit{Ресурс} & \textit{Срок} \\
    Мягк. навыки & Освоение деловой коммуникации с заказчиком & Изучение курса «Основы деловой коммуникации» (по учебному плану) & Успешная сдача зачета & Время & $6$ семестр \\
    Мягк. навыки & Юридическая грамотность & Изучение ФЗ о налоге на профессиональный доход & Регистрация в приложении «Мой налог» & Время & $1$ мес. \\
    Мягк. навыки & Самопрезентация & Упаковка профилей на биржах фриланса & Заполненное портфолио & Время & $2$ нед. \\
    Тв. навыки & Расширение стека & Изучение методов работы с REST API & Разработанный бот & Время, ПК & $1$ мес. \\
    Мягк. навыки & Финансовое планирование & Внедрение системы учета доходов и расходов & Таблица доходов и расходов & Время & $1$ нед. \\
\end{longtblr}

\subsection{Оценка чистого дохода (УСЛУГИ)}

Доход при самозанятости формируется исходя из почасовой ставки и объема отработанного времени. 
Услованая планируемая плата за услугу ($\text{ПУ}$) составляет $700$ руб./час. 
Планируемая нагрузка: $6$ часов в день, $5$ дней в неделю. Среднее количество рабочих дней в месяце принимается равным $22$.
Объем рабочих часов в месяц ($\text{Ч}_{\text{мес}}$):
\begin{equation}
    \text{Ч}_{\text{мес}} = 6 \cdot 22 = 132 \text{ часа.}
\end{equation}

Потенциальный доход ($\text{ПД}$) в месяц составит:
\begin{equation}
    \text{ПД} = \text{ПУ} \cdot \text{Ч}_{\text{мес}} = 700 \cdot 132 = 92400 \text{ руб.}
\end{equation}

Расходы на продвижение ($\text{У}$) на начальном этапе равны $0$ руб.
Обязательный минимальный взнос на пенсионное страхование в 2026 году составляет $5961$ руб./мес.

Расчет чистого потенциального дохода ($\text{ЧД}_{\text{мес}}$) производится в двух вариантах в зависимости от статуса заказчика.
При оказании услуг физическим лицам (налог $4\%$):
\begin{equation}
    \text{ЧД}_{\text{мес}} = \text{ПД} \cdot \frac{100\% - 4\%}{100\%} - 5961 - \text{У} = 92400 \cdot 0.96 - 5961 - 0 = 82743 \text{ руб.}
\end{equation}

При оказании услуг юридическим лицам (налог $6\%$):
\begin{equation}
    \text{ЧД}_{\text{мес}} = \text{ПД} \cdot \frac{100\% - 6\%}{100\%} - 5961 - \text{У} = 92400 \cdot 0.94 - 5961 - 0 = 80895 \text{ руб.}
\end{equation}

Таким образом, чистый потенциальный доход при выбранной нагрузке составит от $80895$ до $82743$ руб. в месяц.